\begin{abstract}
Vision-based tactile sensors (VBTS), such as the Gelsight, capture surface detail through images of a deformable gel, but these images also contain patterns and illumination specific to the sensor hardware. We show that material classifiers exploit this hardware state as a shortcut, compromising recognition after sensor or gel replacement. Across four months of recordings covering 2,711 robot grasps, 20 materials, two GelSight Mini units, and three gels, the state remains relatively stable within each hardware configuration and changes abruptly when the unit or gel is replaced. When materials are associated with particular hardware during training, classifiers mistake materials observed after replacement for those previously recorded on the new hardware. Controlled experiments on public datasets confirm this failure mechanism. We address it using the robot’s own unlabelled contacts to estimate the current hardware state. A training-free correction removes the smooth illumination difference between this running estimate and the training data, while a change detector resets the estimate after hardware replacement. When retraining is possible, adversarial feature alignment uses the same contacts to address remaining hardware differences. Across 55 sessions whose materials were learned on different hardware, correction alone raises ResNet-50 material recognition accuracy from \XpRecipeAcc\ to \XpReadAcc\pct\ (chance 5\pct), performing comparably to the strongest evaluated test-time adaptation baseline. It also improves five domain adaptation methods by 3–10 percentage points. Combined with alignment, accuracy reaches \XpFullAcc\pct\, with most of the alignment gain retained on sessions excluded from the adaptation stream. We release the four-month dataset, a protocol for evaluation across hardware configurations, and code to support tactile recognition under routine hardware changes.
\end{abstract}
